\documentclass[11pt]{article}
\usepackage{parskip}
\usepackage{color}

\title{General cookie recipes}
\author{Nathaniel Zuk}

\begin{document}
\maketitle

\textcolor{green}{(1/6/2018)}

I wanted to see if I could make cookies without any eggs.  How would they turn out?  My expectation was that they would be flat.  I accentuated that (went for crispier cookies) but adding more sugar and cutting down a bit on the flour.

\textbf{Ingredients:}

\begin{enumerate}
\item Butter = 1 cup
\item Dark brown sugar = 1 cup
\item Sugar = 1 cup
\item Milk = 1/4 cup
\item Vanilla = 1 tsp
\item All purpose flour = 2 cups (257 g)
\item Cocoa powder = 1/2 cup (50 g)
\item Baking soda = 1 tsp
\item Salt = 1 tsp
\item Cinnamon = 1 tsp
\end{enumerate}

With oven at 350 F:

\begin{enumerate}
\item 15 min = Cookies turn out thin but still chewy out of the oven.  They harden though as the cool.
\item 12 min = Puffy when I take them out of the oven, and flatten while cooling.  Still malleable but harden a bit when cooling.
\end{enumerate}

With oven at 375 F:

\begin{enumerate}
\item 12 min = Still chewy on the inside and slightly crispier on the outside.  I think they are still chewy after cooling.
\end{enumerate}

\textcolor{blue}{Eggs allow the cookies to hold their shape.  The cookies still rise when they're baking, but they flatten out as the cool.  However, since there are no eggs, there is no risk of salmonella if they are not quite cooked completely.  Shorter cooking times at a higher temperature appear to be better for having them remain chewy after cooling.}

\textcolor{green}{(1/23/2018)}

This time I used less flour.  I also left out the baking soda since the previous cookies I made did not hold air, and I was expecting them to flatten out while cooling.

Ingredients:

\begin{enumerate}
\item 1 cup butter
\item 1 cup dark brown sugar (169 g)
\item 1 cup sugar (222 g)
\item 1/3 cup milk
\item 1 tsp vanilla
\item 1.5 cups flour (225 g)
\item 1/2 cup cocoa powder
\item 1 tsp salt
\item 1 tsp cinnamon
\end{enumerate}

Baked at 375 F for 12 min initially.  The centers of the cookies were still chewy, and I thought that they could be cooked longer, so the next two batches were cooked for 14 min.  

\textcolor{blue}{The resulting cookies were crispy on the outside and chewy inside.  The outer edges of the cookies were a bit to crispy though -- the might be better of the dough melted more evenly or was balled together more tightly before baking.  The butter really melts out, the cookes have holes and cavities where the butter melted away.  They might also look better if the dough is packed more tightly beforehand.

After cooling completely, the next day, the cookies are crunchy and tough.  They are still soft in the middle (a bit), but most of the cookie is crunchy.}

\hrulefill

\textcolor{green}{(1/17/2018)}

I made cranberry cinnamon cookies and tried to make them gluten free using Tina's gluten free flour.  The flour is: 

\begin{enumerate}
\item 1 part white rice flour
\item 1 part tapioca flour
\item 2 parts potato starch
\item 6 parts brown rice flour
\end{enumerate}

Then I added 1/2 tsp of xantham gum.  The rest of the ingredients were the same as before:

\textbf{Ingredients:}

\begin{enumerate}
\item Butter = 1 cup (2 sticks)
\item Sugar = 3/4 cup
\item Brown sugar = 3/4 cup
\item Vanilla = 1 tsp
\item Eggs = 2
\item Baking soda = 1 tsp
\item Salt = 1 tsp
\item Cranberries ~ 1 cup \textcolor{red}{(I didn't measure this, but probably added more)}
\item Cinnamon = 1/2 tsp
\item Gluten free flour mix = 2 1/4 cups
\item Xantham gum (see above)
\end{enumerate}

Baked at 375 F for 15 min.  

\textcolor{blue}{The initial batch turned out too crumbly.  I had used up about 2/3 of the batter in the first batch.  The cookies were pretty large and spread, which may have been because I used a lot of batter for each cookies (I didn't measure the amount, but probably 1-2 inches in diameter).  

I split the next batch in two.  For one I added 1 tsp of egg whites.  For the other I added 1/2 tsp xantham gum.  I also made the cookies smaller (~1 inch in diameter).  I cooked these for 12 min since they looked cooked after that time.  The resulting cookes held together better than the last batch, and the xantham gum ones held together better than the egg white ones.  

Tina also recommended using less butter to keep them from spreading out so much.}

\end{document}